\documentclass[11pt]{article}
\usepackage[T1]{fontenc}
\usepackage{textcomp}
\usepackage[margin=0.75in]{geometry}
\usepackage{amsmath,amssymb,amsthm}
\usepackage{array,booktabs}
\usepackage{hyperref}
\setlength{\parindent}{0pt}
\newtheorem{theorem}{Theorem}
\theoremstyle{definition}
\newtheorem{definition}{Definition}
\title{Flow Proof Artifact: prop-02-copy-segment}
\author{Generated by \texttt{flow doc proof}}
\date{Flow Proof Book}
\begin{document}
\maketitle
To place at a given point a straight line equal to a given straight line.\\[0.5em]
\textbf{Source.} euclid — Elements, Book I, Proposition 2\\[0.5em]
\bigskip
\noindent{\large \textbf{Axiom 1.} \textit{To place at a given point a straight line equal to a given straight line} \hfill the Euclidean plane \cdot Euclid Book I \cdot Proposition 2: a segment equal to a given segment}\par
\smallskip\noindent\textit{Coordinate: the Euclidean plane \cdot Euclid Book I \cdot Proposition 2: a segment equal to a given segment}\par
\smallskip\noindent\textit{Source: euclid — Elements, Book I, Proposition 2}\par
\medskip
\noindent\textbf{Goal.} To place at a given point a straight line equal to a given straight line\par
\noindent$\displaystyle \text{segment at A equals BC}$\par
\medskip
\noindent
\begin{tabular}{@{} >{\raggedright\arraybackslash}p{0.06\textwidth} >{\raggedright\arraybackslash}p{0.47\textwidth} >{\raggedleft\arraybackslash}p{0.41\textwidth} @{}}
\textbf{\#} & \textbf{Proof} & \textbf{Mathematics} \\
\midrule
\textbf{1.} & We accept proposition 2: a segment equal to a given segment for Book I of the Elements in the Euclidean plane without proof — an ontological commitment, not a lemma. & \\[0.45em]
\textbf{2.} & Let L = line\_through\_A\_and\_D. & \\[0.45em]
\textbf{3.} & Let equilateral\_ADE = triangle\_on\_AD\_by\_Proposition\_1. & \\[0.45em]
\textbf{4.} & Let circle\_center\_D\_radius\_BC = circle\_at\_D\_with\_radius\_BC. & \\[0.45em]
\textbf{5.} & Let circle\_center\_E\_radius\_DG = circle\_at\_E\_with\_radius\_DG. & \\[0.45em]
\textbf{6.} & Let F = second\_intersection\_of\_circles\_at\_D\_and\_E. & \\[0.45em]
\textbf{7.} & We invoke the derived fact: Postulate 3: a circle can be drawn with any center and radius. & \\[0.45em]
\textbf{8.} & We invoke the derived fact: Postulate 1: a straight line can be drawn from any point to any point. & \\[0.45em]
\textbf{9.} & From step 2, step 3, step 4, step 5, step 6, step 7, and step 8, we can deduce that segment AF equals segment BC. & $\displaystyle \text{segment AF} = \text{segment BC}$ \\[0.45em]
\textbf{10.} & This holds immediately by the stated axiom: segment at A equals BC. Hence proven. & $\displaystyle \text{segment at A equals BC}$ \\[0.45em]
\bottomrule
\end{tabular}
\label{thm:Geometry-Euclid Book I-proposition_2_a_segment_equal_to_a_given_segment}
\par\medskip\hrule
\end{document}
